\documentclass[11pt]{article}
\usepackage[margin=1in]{geometry}
\usepackage{amsmath,amssymb,amsthm}
\usepackage{booktabs}
\usepackage{hyperref}
\usepackage{mathrsfs}

\newtheorem{theorem}{Theorem}
\newtheorem{lemma}[theorem]{Lemma}
\newtheorem{proposition}[theorem]{Proposition}
\newtheorem{conjecture}[theorem]{Conjecture}
\newtheorem{definition}[theorem]{Definition}
\newtheorem{corollary}[theorem]{Corollary}

\title{Prime Geometry III:\\
A Generalized Bertrand Principle for Prime Constellations}
\author{Allen Proxmire}
\date{April 2026}

\begin{document}
\maketitle

\begin{abstract}
In PG~II we introduced the Twin-Prime Bertrand Postulate
$(\mathrm{TPB})$: every interval $(x,2x]$ with $x\ge 11$ contains a
twin prime pair. Here we show that the same dyadic inequality
appears to hold, with the same universal factor $2$, for every
admissible pair-constellation. We formulate the Generalized Bertrand
Principle $(\mathrm{GBP})$:
\[
\pi_\mathcal{C}(2x)-\pi_\mathcal{C}(x)\ge 1\qquad(x\ge P^{*}_\mathcal{C})
\]
for every admissible constellation $\mathcal{C}=\{0,g\}$ with
Hardy--Littlewood singular series $\mathfrak{S}(\mathcal{C})>0$. We
give a generalized $2P$-beats lemma linking $(\mathrm{GBP})$ to the
Prime Triangle angle geometry, prove the equivalence between the
dyadic and ratio formulations, and extract a data-driven envelope
\[
G^{\mathcal{C}}_j \;<\; 0.171\,(\log P^{\mathcal{C}}_j)^{3.22}
\]
that fits all three constellations with $R^2\ge 0.98$. Empirical verification uses
the twin ($g=2$), cousin ($g=4$) and sexy-prime ($g=6$) sequences up
to $10^{10}$: $27\,412\,679$ twins, $27\,409\,999$ cousins and
$54\,818\,296$ sexy pairs, with $P^{*}_{\{0,2\}}=11$,
$P^{*}_{\{0,4\}}=7$, and $P^{*}_{\{0,6\}}=5$. We discuss the
asymptotic interpretation, the drift of the twin-gap exponent
$\beta$ toward the Hardy--Littlewood value $2$, and the apparent
universality of the overshoot law
$\Omega(T)\sim(\log\log T)^{\sim 1}$.
\end{abstract}

\section{Introduction}

PG~II isolated a single inequality~---~the Twin-Prime Bertrand
Postulate~---~as a natural target in twin-prime distribution:
\[
\pi_2(2x)-\pi_2(x)\ge 1\qquad(x\ge 11).
\]
It was verified there for $x\le 10^9$; we have since extended the
verification to $x\le 10^{10}$, still without exception beyond the
trivial small-$x$ case $T_k=5$. The inequality has an equivalent
geometric reading: every angle-record of the Prime Triangle sequence
with $p_n\ge 3$ is realized by a twin prime.

The factor $2$ in $(\mathrm{TPB})$ is a Bertrand factor~---~it matches
the dyadic interval $(x,2x]$~---~but until now we had not asked
whether it was specific to twins. In this note we investigate the
behavior of other prime constellations (the cousin primes,
$\{0,4\}$, and the sexy primes, $\{0,6\}$) and find that the same
bound holds:

\begin{center}
\emph{For every admissible pair-constellation, the dyadic inequality
$\pi_\mathcal{C}(2x)-\pi_\mathcal{C}(x)\ge 1$ holds for all
sufficiently large $x$, with a threshold $P^{*}_\mathcal{C}$ that is
very small in every case examined.}
\end{center}

We call this the Generalized Bertrand Principle $(\mathrm{GBP})$ and
establish its equivalence with an elementary ratio form and, via a
generalized $2P$-beats lemma, with a family of angle-record
statements in the Prime Triangle sense. We verify the principle for
all three constellations to $10^{10}$, and combine the data into a
uniform empirical envelope for the extreme gaps.

Throughout, $\mathbb{P}$ denotes the primes. The Hardy--Littlewood
prime-tuple conjecture predicts, for any admissible constellation
$\mathcal{C}$, an asymptotic twin-like density
$\pi_\mathcal{C}(x)\sim 2\,\mathfrak{S}(\mathcal{C})\,x/(\log x)^2$,
with $\mathfrak{S}(\mathcal{C})$ the singular series. For the three
constellations studied here, $\mathfrak{S}(\{0,2\})=
\mathfrak{S}(\{0,4\})=C_2\approx 0.6602$ and
$\mathfrak{S}(\{0,6\})=2C_2$ (see Section~\ref{sec:asymp}).

\section{Admissible Constellations and $(\mathrm{GBP})$}

\begin{definition}[Admissible pair-constellation]
A \emph{pair-constellation} is a set $\mathcal{C}=\{0,g\}$ with
$g\ge 2$ even. It is \emph{admissible} if for every prime $q$ the
reduction $\mathcal{C}\pmod q$ does not cover all residue classes.
Equivalently, the singular series
\[
\mathfrak{S}(\mathcal{C})
=\prod_{q\text{ prime}}\frac{1-\nu_\mathcal{C}(q)/q}{(1-1/q)^{|\mathcal{C}|}}
\]
is positive, where $\nu_\mathcal{C}(q)$ is the number of residues
modulo $q$ occupied by $\mathcal{C}$.
\end{definition}

For every admissible $\mathcal{C}=\{0,g\}$, let
\[
P^{\mathcal{C}}_1<P^{\mathcal{C}}_2<\cdots
\]
be the sequence of leading members: $P^{\mathcal{C}}_j\in\mathbb{P}$
with $P^{\mathcal{C}}_j+g\in\mathbb{P}$.

\begin{definition}[$\mathcal{C}$-counting function]
$\pi_\mathcal{C}(x)=\#\{p\le x:p\in\mathbb{P},\,p+g\in\mathbb{P}\}$.
\end{definition}

\begin{conjecture}[Generalized Bertrand Principle, $(\mathrm{GBP})$]
\label{conj:GBP}
For every admissible pair-constellation $\mathcal{C}=\{0,g\}$ there
exists a threshold $P^{*}_\mathcal{C}$ such that
\[
\pi_\mathcal{C}(2x)-\pi_\mathcal{C}(x)\ge 1\qquad(x\ge P^{*}_\mathcal{C}).
\]
Equivalently, $P^{\mathcal{C}}_{j+1}<2\,P^{\mathcal{C}}_j$ for every
$j$ with $P^{\mathcal{C}}_j\ge P^{*}_\mathcal{C}$.
\end{conjecture}

\begin{proposition}[Dyadic and ratio forms]\label{prop:dyadic}
For each admissible $\mathcal{C}$ with fixed threshold
$P^{*}_\mathcal{C}$, the two formulations above are equivalent.
\end{proposition}

\begin{proof}
The proof is identical to Proposition~2.1 of PG~II. If
$P^{\mathcal{C}}_{j+1}<2P^{\mathcal{C}}_j$ whenever
$P^{\mathcal{C}}_j\ge P^{*}_\mathcal{C}$, fix $x\ge P^{*}_\mathcal{C}$
and let $P^{\mathcal{C}}_j$ be the largest term
$\le x$; then $P^{\mathcal{C}}_{j+1}<2P^{\mathcal{C}}_j\le 2x$, giving
at least one $\mathcal{C}$-pair in $(x,2x]$. Conversely, if the
dyadic inequality holds and $P^{\mathcal{C}}_j\ge P^{*}_\mathcal{C}$,
the dyadic form with $x=P^{\mathcal{C}}_j$ produces
$P^{\mathcal{C}}_{j+1}\in(P^{\mathcal{C}}_j,2P^{\mathcal{C}}_j]$; strict
inequality follows because $2P^{\mathcal{C}}_j$ is even and
$P^{\mathcal{C}}_{j+1}$ is odd (for $g>0$ even, both members of the
pair are odd).
\end{proof}

Observe that $(\mathrm{TPB})$ is the special case
$\mathcal{C}=\{0,2\}$ with $P^{*}_{\{0,2\}}=11$.

\section{A Generalized $2P$-Beats Lemma}\label{sec:2Pgen}

Recall from PG~I the Prime Triangle angle of a pair $(p,q)$ with
$p<q$:
\[
\alpha(p,q)=\arctan(p/q),\qquad \rho(p,q)=p/q.
\]
For a fixed constellation $\mathcal{C}=\{0,g\}$ and a member
$P^{\mathcal{C}}_j$, the pair $(P^{\mathcal{C}}_j,P^{\mathcal{C}}_j+g)$
has angle $\alpha^{\mathcal{C}}_j=\arctan(P^{\mathcal{C}}_j/(P^{\mathcal{C}}_j+g))$.

\begin{lemma}[Generalized $2P$-beats]\label{lem:gen2P}
Let $\mathcal{C}=\{0,g\}$ and $\mathcal{D}=\{0,h\}$ with $g\ne h$, and
let $(P,P+g)$, $(Q,Q+h)$ be prime pairs for $\mathcal{C}$ and
$\mathcal{D}$ respectively. Then
\[
\rho(P,P+g)>\rho(Q,Q+h)\iff h\,P>g\,Q.
\]
In particular:
\begin{enumerate}
\item[(i)] For $h=2,\,g=4$: a cousin pair $(P,P+4)$ beats a twin pair
$(Q,Q+2)$ in angle iff $P>2Q$ (the $2P$-beats rule of PG~II).
\item[(ii)] For $h=g$: both pairs belong to the same constellation
$\mathcal{C}$, and $\rho(P,P+g)>\rho(Q,Q+g)\iff P>Q$. The angle
sequence within a fixed $\mathcal{C}$ is monotonic in the leading
prime.
\item[(iii)] For a constellation $\mathcal{C}=\{0,g\}$ and its
``doubled companion'' $\mathcal{C}'=\{0,2g\}$, a $\mathcal{C}'$-pair
at $P$ beats a $\mathcal{C}$-pair at $Q$ iff $P>2Q$. The factor $2$
is universal for the $\mathcal{C}\leftrightarrow\mathcal{C}'$
comparison, independent of $g$.
\end{enumerate}
\end{lemma}

\begin{proof}
$\frac{P}{P+g}>\frac{Q}{Q+h}\iff P(Q+h)>Q(P+g)\iff hP>gQ$. The
specializations follow.
\end{proof}

Part~(iii) gives the geometric explanation for the universal factor
$2$ in $(\mathrm{GBP})$: for every admissible $\mathcal{C}=\{0,g\}$,
the \emph{doubled-gap} companion $\mathcal{C}'=\{0,2g\}$ is also
admissible (its singular series is positive), and the two
constellations are related by the same $2P$-beats threshold. The
dyadic factor of $2$ in the Bertrand inequality is the natural
density counterpart of the same factor in the angle comparison.

\section{The Generalized Angle-Record Theorem}\label{sec:gen-angle}

For twins, PG~II's Angle-Record Theorem states that every record-setting
$\alpha_n$ among consecutive prime pairs $(p_n,p_{n+1})$ with
$p_n\ge 3$ is a twin pair, and that this is equivalent to
$(\mathrm{TPB})$.

To extend this to a general $\mathcal{C}$ we need two things. First,
a \emph{domain} on which $\mathcal{C}$-pairs can plausibly dominate
the angle-records: if $\mathcal{C}$ has gap $g>2$, twins always beat
$\mathcal{C}$-pairs of comparable leading prime
(since $\rho(P,P+g)<\rho(P,P+2)$), and twin records preempt
$\mathcal{C}$-records in the unrestricted sequence.

Fix $\mathcal{C}=\{0,g\}$ and consider only the pairs
$(p,p+g')$ with $g'\ge g$ and both entries prime~---~call these
\emph{$\mathcal{C}$-admissible pairs}. Within this restricted
universe, the $2P$-beats threshold for a competing pair of strictly
larger gap $g'>g$ is $P>(g'/g)Q\ge(g+2)/g\cdot Q$, so the
\emph{binding} competing threshold is the companion
$\mathcal{C}'=\{0,g+2\}$ with factor $(g+2)/g$. Only in the special
case $g=2$ (twins, with companion cousins) is this binding factor
equal to~$2$.

\begin{definition}[$\mathcal{C}$-angle-record]
A $\mathcal{C}$-admissible pair $(P^*,P^*+g^*)$ with $g^*\ge g$ is a
\emph{$\mathcal{C}$-angle-record} if
$\rho(P^*,P^*+g^*)>\rho(P,P+g')$ for every earlier
$\mathcal{C}$-admissible pair $(P,P+g')$.
\end{definition}

\begin{theorem}[Generalized Angle-Record Theorem]
\label{thm:gen-angle}
For each admissible $\mathcal{C}=\{0,g\}$, the following are
equivalent:
\begin{enumerate}
\item[(i)] $P^{\mathcal{C}}_{j+1}<\frac{g+2}{g}\,P^{\mathcal{C}}_j$
for every $j$ with $P^{\mathcal{C}}_j\ge P^{\star}_\mathcal{C}$ (for
some effective $P^{\star}_\mathcal{C}$ depending on the gap).
\item[(ii)] Every $\mathcal{C}$-angle-record with
$P^*\ge P^{\star}_\mathcal{C}$ is a pair of constellation $\mathcal{C}$.
\end{enumerate}
For $g=2$, $(g+2)/g=2$ and $(i)$ reduces to $(\mathrm{TPB})$;
Theorem~\ref{thm:gen-angle} is then PG~II's Angle-Record Theorem.
For $g=4$ the factor is $1.5$; for $g=6$ the factor is $4/3$; and so
on.
\end{theorem}

\begin{proof}[Proof sketch]
The argument is analogous to the twin case. (i)$\Rightarrow$(ii):
suppose (i) holds. Between two consecutive $\mathcal{C}$-members
$P^{\mathcal{C}}_j$ and $P^{\mathcal{C}}_{j+1}$, a competing
$\mathcal{C}$-admissible pair $(P,P+g')$ with $g'>g$ beats the
$\mathcal{C}$-record at $P^{\mathcal{C}}_j$ iff $P>(g'/g)P^{\mathcal{C}}_j$;
the binding case is $g'=g+2$, requiring
$P>((g+2)/g)P^{\mathcal{C}}_j$. Under (i),
$P^{\mathcal{C}}_{j+1}<((g+2)/g)P^{\mathcal{C}}_j$, so no competing
pair in the window can exceed the threshold, and the next
$\mathcal{C}$-admissible record is either a larger-$P$ pair of gap
$g$ (i.e., $P^{\mathcal{C}}_{j+1}$) or none until then; meanwhile
$P^{\mathcal{C}}_{j+1}$ does set a new record because
$\rho(P^{\mathcal{C}}_{j+1},P^{\mathcal{C}}_{j+1}+g)>
\rho(P^{\mathcal{C}}_j,P^{\mathcal{C}}_j+g)$.

(ii)$\Rightarrow$(i): contrapositive. If (i) fails for some $j$,
there exist $(P,P+g')$ with $g'>g$, both primes, and
$P^{\mathcal{C}}_j<P\le P^{\mathcal{C}}_{j+1}$ with
$P>((g+2)/g)P^{\mathcal{C}}_j$. By
Lemma~\ref{lem:gen2P}, this pair beats the previous
$\mathcal{C}$-record, so a non-$\mathcal{C}$-record appears. The
three smallest anomalies (very small $P$) are handled by direct
inspection, producing the effective $P^{\star}_\mathcal{C}$.
\end{proof}

\section{Empirical Verification up to $10^{10}$}\label{sec:empirical}

All data in this section derive from a single segmented (odd-only)
sieve of the primes up to $10^{10}$, followed by a streaming
extraction of all pairs $(p,p+g)$ with both primes for
$g\in\{2,4,6\}$. Cross-boundary pairs are handled by a three-entry
buffer. The sieve confirms $455\,052\,511$ primes below $10^{10}$
and the following constellation counts:

\begin{center}
\begin{tabular}{lrr}
\toprule
constellation $\mathcal{C}$ & gap $g$ & count up to $10^{10}$ \\
\midrule
twin & 2 & $27\,412\,679$ \\
cousin & 4 & $27\,409\,999$ \\
sexy & 6 & $54\,818\,296$ \\
\bottomrule
\end{tabular}
\end{center}

\subsection{Ratio envelopes}

For each constellation we tabulate $\sup r^{\mathcal{C}}_k$ on tails
$P^{\mathcal{C}}_j>X$:

\begin{center}
\begin{tabular}{lrrr}
\toprule
scope & twin sup $r$ & cousin sup $r$ & sexy sup $r$ \\
\midrule
all & $2.200000$ & $2.333333$ & $1.571429$ \\
$P_j>100$ & $1.280374$ & $1.283465$ & $1.224299$ \\
$P_j>10^3$ & $1.081880$ & $1.109790$ & $1.051647$ \\
$P_j>10^6$ & $1.000711$ & $1.000865$ & $1.000695$ \\
$P_j>10^9$ & $1.000004$ & $1.000004$ & $1.000002$ \\
\bottomrule
\end{tabular}
\end{center}

The three envelopes converge to within $10^{-5}$ of $1$ by
$P_j>10^9$, with comparable decay rates.

\subsection{Bertrand threshold $P^{*}_\mathcal{C}$}

The top ratios in each constellation (ranks~1--5 shown):

\begin{center}
\begin{tabular}{lrrrr}
\toprule
constellation & $P_j$ & $P_{j+1}$ & $r_j$ & gap \\
\midrule
twin & $5$ & $11$ & $2.2000$ & $6$ \\
twin & $17$ & $29$ & $1.7059$ & $12$ \\
twin & $3$ & $5$ & $1.6667$ & $2$ \\
twin & $11$ & $17$ & $1.5455$ & $6$ \\
twin & $41$ & $59$ & $1.4390$ & $18$ \\
\midrule
cousin & $3$ & $7$ & $2.3333$ & $4$ \\
cousin & $19$ & $37$ & $1.9474$ & $18$ \\
cousin & $7$ & $13$ & $1.8571$ & $6$ \\
cousin & $43$ & $67$ & $1.5581$ & $24$ \\
cousin & $13$ & $19$ & $1.4615$ & $6$ \\
\midrule
sexy & $7$ & $11$ & $1.5714$ & $4$ \\
sexy & $5$ & $7$ & $1.4000$ & $2$ \\
sexy & $17$ & $23$ & $1.3529$ & $6$ \\
sexy & $23$ & $31$ & $1.3478$ & $8$ \\
sexy & $13$ & $17$ & $1.3077$ & $4$ \\
\bottomrule
\end{tabular}
\end{center}

Violations of the strong $r<2$ bound:

\begin{center}
\begin{tabular}{lrr}
\toprule
constellation & \# violations of $r_j<2$ & $P^{*}_\mathcal{C}$ \\
\midrule
twin & $1$ (only $(5,7)\!\to\!(11,13)$) & $11$ \\
cousin & $1$ (only $(3,7)\!\to\!(7,11)$) & $7$ \\
sexy & $0$ & $5$ \\
\bottomrule
\end{tabular}
\end{center}

This verifies $(\mathrm{GBP})$ with the stated thresholds for all
three constellations up to $10^{10}$. In particular, GBP holds with
the \emph{same} factor $2$ across constellations; the user's
suggestion of larger factors ($3$ for cousins, $4$ for sexy) is
true but unnecessarily loose~---~the $r<2$ bound captures the
structure exactly.

\subsection{Per-decade twin overshoot}

The twin overshoot
$\Omega(T)=\max G_k/(\log T_k^{\star})^2$ per decade (stable regime
$10^3\!\le\!T_k\!<\!10^9$):

\begin{center}
\begin{tabular}{lrrrr}
\toprule
decade & max $G_k$ & at $T_k$ & $(\log T)^2$ & $\Omega$ \\
\midrule
$10^3$ & $210$ & $5\,879$ & $75.33$ & $2.788$ \\
$10^4$ & $630$ & $62\,297$ & $121.87$ & $5.169$ \\
$10^5$ & $1\,452$ & $850\,349$ & $186.42$ & $7.789$ \\
$10^6$ & $1\,722$ & $9\,923\,987$ & $259.55$ & $6.635$ \\
$10^7$ & $2\,868$ & $96\,894\,041$ & $338.16$ & $8.481$ \\
$10^8$ & $4\,770$ & $698\,542\,487$ & $414.71$ & $11.502$ \\
$10^9$ & $6\,030$ & $4\,289\,385\,521$ & $491.93$ & $12.258$ \\
\bottomrule
\end{tabular}
\end{center}

\section{A Refined Envelope Conjecture}\label{sec:envelope}

For each constellation we compute $\sup(r_j-1)$ over tails
$P^{\mathcal{C}}_j>T_\min$ at
$T_\min\in\{10^2,10^3,10^5,10^7,10^9\}$:

\begin{center}
\begin{tabular}{lrrr}
\toprule
$T_\min$ & twin $\sup(r-1)$ & cousin $\sup(r-1)$ & sexy $\sup(r-1)$ \\
\midrule
$10^2$ & $2.80\times10^{-1}$ & $2.83\times10^{-1}$ & $2.24\times10^{-1}$ \\
$10^3$ & $8.19\times10^{-2}$ & $1.10\times10^{-1}$ & $5.16\times10^{-2}$ \\
$10^5$ & $5.10\times10^{-3}$ & $4.33\times10^{-3}$ & $3.34\times10^{-3}$ \\
$10^7$ & $1.43\times10^{-4}$ & $1.49\times10^{-4}$ & $9.87\times10^{-5}$ \\
$10^9$ & $3.86\times10^{-6}$ & $4.02\times10^{-6}$ & $2.26\times10^{-6}$ \\
\bottomrule
\end{tabular}
\end{center}

The three envelopes decay at essentially the same rate. A weighted
least-squares fit of the log-linear form
\[
\log\!\big((r_j-1)T_j\big)
=\log C_\mathcal{C}+\delta_\mathcal{C}\log\log T_j
\]
on the five $T_\min$ points yields:

\begin{center}
\begin{tabular}{lrrr}
\toprule
constellation & $\delta_\mathcal{C}$ & $C_\mathcal{C}$ & $R^2$ \\
\midrule
twin & $3.2884$ & $0.1645$ & $0.9970$ \\
cousin & $3.2209$ & $0.2013$ & $0.9957$ \\
sexy & $3.1538$ & $0.1507$ & $0.9931$ \\
\midrule
pooled (uniform fit) & $3.2210$ & $0.1709$ & $0.9824$ \\
\bottomrule
\end{tabular}
\end{center}

The fitted exponents agree across constellations within $4.2\%$
(range $[3.15,3.29]$), strongly suggesting a common asymptotic
$\delta$. The constants $C_\mathcal{C}$ vary by roughly $29\%$
(range $[0.15,0.20]$); this larger spread is driven by the cousin
envelope, whose leading offenders at small $P_j$ (dominated by the
pair $(3,7)\to(7,11)$) inflate its fitted prefactor relative to the
asymptotic regime.

The pooled fit~---~obtained by concatenating all fifteen data
points and fitting a single $(C,\delta)$ pair~---~yields
\begin{equation}\label{eq:env-uniform}
r_j-1 \;\le\; 0.171\,\frac{(\log P_j)^{3.22}}{P_j},
\qquad
G^{\mathcal{C}}_j \;<\; 0.171\,(\log P^{\mathcal{C}}_j)^{3.22}
\qquad(R^2=0.982).
\end{equation}

\begin{conjecture}[Refined envelope, data-driven]\label{conj:env}
For every admissible pair-constellation $\mathcal{C}=\{0,g\}$ there
exist constants $C_\mathcal{C},\delta_\mathcal{C}$ with
$\delta_\mathcal{C}\to\delta^*$ as $g$ varies over admissible gaps,
for some universal $\delta^*\in[3.15,3.30]$, such that
\[
G^\mathcal{C}_j \;=\; P^{\mathcal{C}}_{j+1}-P^{\mathcal{C}}_j
\;<\; C_\mathcal{C}\,(\log P^{\mathcal{C}}_j)^{\delta_\mathcal{C}}
\qquad(P^{\mathcal{C}}_j\ge P^{*}_\mathcal{C}).
\]
The pooled uniform envelope~\eqref{eq:env-uniform} fits all three
examined constellations with $R^2\ge 0.98$.
\end{conjecture}

This is strictly stronger than $(\mathrm{GBP})$, which is recovered
by the coarse bound $G^\mathcal{C}_j/P^{\mathcal{C}}_j\to 0$. It is
compatible with the Hardy--Littlewood prediction for the
\emph{typical} gap $G^\mathcal{C}_j\sim(\log P)^2/(2\mathfrak{S}(\mathcal{C}))$:
the exponent $\delta\approx 3.2>2$ captures the \emph{extreme} gap,
which exceeds the typical gap by an overshoot factor that itself
grows slowly with $P$.

\section{Asymptotic Interpretation}\label{sec:asymp}

\subsection{Exponent drift toward $\beta=2$}

In PG~II the twin-gap power law $G_k\approx A(\log T_k)^\beta$ was
fitted on primes up to $10^9$, yielding $\beta\approx 1.866$ stable
across sub-ranges. Extension to $10^{10}$ shifts the fits upward:

\begin{center}
\begin{tabular}{lrr}
\toprule
fit range $T_k>$ & $\beta$ ($10^9$ data) & $\beta$ ($10^{10}$ data) \\
\midrule
$10^3$ & $1.8633$ & $1.8865$ \\
$10^5$ & $1.8659$ & $1.8872$ \\
$10^7$ & $1.8664$ & $1.8889$ \\
$10^9$ & ---      & $1.8957$ \\
\bottomrule
\end{tabular}
\end{center}

The drift is monotone; each fit band grew $\Delta\beta\approx+0.02$
across the new decade. Under the interpretation that the true
asymptotic exponent is $\beta=2$ (Hardy--Littlewood), our observations
are consistent with a slow finite-size convergence whose pace is
controlled by the overshoot factor
$\Omega$.

\subsection{Overshoot growth law}

Fits of the per-decade twin overshoot
$\Omega(T)$ on the seven stable points yield:
\begin{center}
\begin{tabular}{lll}
\toprule
model & form & $R^2$ \\
\midrule
(a) & $\Omega=A(\log T)^b$, $A=0.147,\,b=1.43$ & $0.9133$ \\
(b) & $\log\Omega=a+b(\log\log T)^\alpha$, $\alpha=0.69$ & $0.9135$ \\
(c) & $\Omega=D(\log T)^\gamma$, $\gamma=0.10$ & $0.8964$ \\
(d) & $\Omega=A+B\log\log T$, $A=-17.48,B=9.32$ & $0.8959$ \\
\bottomrule
\end{tabular}
\end{center}

Models (a) and (b) are essentially tied, indicating
$\Omega\sim(\log T)^{1.43}$ or equivalently
$\log\Omega\sim(\log\log T)^{0.7}$. This is a
\emph{Cram\'er--Granville}-type growth: Cram\'er conjectured
$p_{n+1}-p_n=O((\log p)^2)$ for ordinary primes, refined by Granville
to include a $\log\log p$ factor. For twin-prime gaps we see an
analogous correction with an exponent $\alpha\approx 0.7$ in the
$\log\log$ factor, slightly sub-logarithmic but firmly positive.

\subsection{Why factor $2$ is natural and universal}

The factor $2$ in $(\mathrm{GBP})$ has a clean interpretation: it
is the dyadic Bertrand factor. For ordinary primes, Bertrand's
postulate ensures a prime in every $(x,2x]$; for any admissible
$\mathcal{C}$, the expected count of $\mathcal{C}$-pairs in $(x,2x]$
under Hardy--Littlewood is
\[
\pi_\mathcal{C}(2x)-\pi_\mathcal{C}(x)
\sim 2\mathfrak{S}(\mathcal{C})\,\frac{x}{(\log x)^2},
\]
which diverges to infinity. Unconditionally we do not know that this
count is $\ge 1$ for every sufficiently large $x$, but the data
indicate that the inequality becomes binding extremely early
($P^{*}_\mathcal{C}$ at most the first few pair members in every
case we examined) and is maintained thereafter with orders of
magnitude of slack.

Lemma~\ref{lem:gen2P}(iii) provides the geometric explanation: the
factor $2$ is exactly the $2P$-beats threshold for the
doubled-companion relation
$\mathcal{C}\leftrightarrow\mathcal{C}'=\{0,2g\}$. In the angle
picture, moving from a $\mathcal{C}$-member at $Q$ to a
$\mathcal{C}'$-member at $P$ requires $P>2Q$ to beat the
$\mathcal{C}$-angle; this threshold is universal in $g$. The dyadic
factor of $2$ in Bertrand's inequality is the density counterpart.

\subsection{Why violations occur only at tiny $P$}

The near-misses of the bound $r<2$ all lie at
$P_j\le 881$ across all three constellations; past this, no ratio in
our data exceeds $1.1$. The reason is purely statistical: at small
$P$, the singular series effectively constrains the positions of
constellation pairs, and a small neighborhood of primes may fail to
contain both $p$ and $p+g$. Once $P$ is large enough that the
expected count in $(P,2P]$ exceeds $1$ by many standard deviations,
fluctuations cannot drive it to zero. The Hardy--Littlewood estimate
gives the expected count
$\sim 2\mathfrak{S}(\mathcal{C})P/(\log P)^2$, which equals $10$ by
$P\sim 10^4$ (twin/cousin) or $P\sim 10^3$ (sexy); a factor of $10$
is already too large for statistical fluctuations to explain a
complete absence.

\section{Conclusion}

We have identified the Generalized Bertrand Principle~---~a
conjectural dyadic inequality $\pi_\mathcal{C}(2x)-\pi_\mathcal{C}(x)
\ge 1$ with a universal factor of $2$~---~as the natural
extension of PG~II's twin-prime Bertrand postulate to every
admissible pair-constellation. The principle is equivalent to the
simple ratio condition $P^{\mathcal{C}}_{j+1}<2P^{\mathcal{C}}_j$ and
is supported geometrically by a generalized $2P$-beats lemma
expressed in the Prime Triangle angle. We have verified the
principle in three constellations (twins, cousins, sexy primes) up
to $P=10^{10}$, giving $P^{*}_{\{0,2\}}=11$, $P^{*}_{\{0,4\}}=7$,
$P^{*}_{\{0,6\}}=5$.

Three observations flow from the computation:
\begin{enumerate}
\item The cousin-prime count at $10^{10}$ ($27\,409\,999$) is
within $0.01\%$ of the twin-prime count ($27\,412\,679$), confirming
the Hardy--Littlewood prediction of identical densities for
$\{0,2\}$ and $\{0,4\}$.
\item The ratio envelopes tighten at essentially the same rate
across constellations. Fitted exponents $\delta_\mathcal{C}$ agree
within $4.2\%$ across the three constellations
(range $[3.15,3.29]$); a pooled uniform fit gives
$G^{\mathcal{C}}_j<0.171(\log P)^{3.22}$ with $R^2=0.982$.
\item The twin-gap exponent $\beta$ drifts upward from
$1.866$ (at $10^9$) to $1.896$ (at $10^{10}$ restricted to
$T_k>10^9$), consistent with the asymptotic value $\beta=2$ predicted
by Hardy--Littlewood; concurrently the overshoot factor $\Omega(T)$
grows roughly like $(\log T)^{1.4}$ or $(\log\log T)^{0.7}$.
\end{enumerate}

Future work. A conditional proof of $(\mathrm{GBP})$ under
Hardy--Littlewood is immediate along the same lines as
Theorem~5.1 of PG~II. An unconditional proof would require a
short-interval result
$\pi_\mathcal{C}(x+y)-\pi_\mathcal{C}(x)>0$ at dyadic scale $y=x$,
which is strictly beyond current sieve-theoretic capabilities
(bounded-gaps results of Zhang, Maynard, and Polymath8 give such
short-interval existence only at some scales, not all). We expect
that the Maynard multi-dimensional sieve, adapted to dyadic
windows, is the natural apparatus for an attack on GBP.

The apparent universality of the envelope across constellations, and
the unified overshoot law, suggest that the Generalized Bertrand
Principle is part of a larger picture: HL-admissible constellations
behave as coordinated copies of the twin-prime sequence, scaled by
their singular-series constants, and subject to a common growth law
for their extreme gaps. Quantifying this universality
conjecturally~---~in particular, obtaining a closed form for the
asymptotic overshoot rate~---~is the natural next step.

\section*{Data availability}

All computations reproduce from
\texttt{scripts/pg\_structural\_1e10.py} and
\texttt{scripts/pg\_extend\_1e10.py}. Saved datasets
(\texttt{twins\_1e10.npy}, \texttt{cousins\_1e10.npy},
\texttt{sexy\_1e10.npy}, and associated CSVs) are stored under the
\texttt{data/} directory.

\section*{Acknowledgements}

This paper is the third in a series. PG~I introduced the Prime
Triangle construction and derived geometric quantities; PG~II
established the Twin-Prime Bertrand Postulate and the Angle-Record
Theorem; the present note generalizes those results to all
admissible pair-constellations.

\end{document}
